\documentclass[11pt]{amsart}

\usepackage[T1]{fontenc}
\usepackage{lmodern}
\usepackage{microtype}
\usepackage{amsmath,amssymb}
\usepackage[hidelinks]{hyperref}

\newtheorem{theorem}{Theorem}
\newtheorem{lemma}[theorem]{Lemma}

\newcommand{\ii}{\mathrm{i}}
\newcommand{\rev}{\mathrm{R}}

\title[Counterexamples to a compression conjecture]
{Counterexamples to the $(v/2)$-compression conjecture for periodic Golay pairs}

\author{Owen Shen}
\address{Operations Research Center, Massachusetts Institute of Technology,
Cambridge, MA 02139, USA}
\email{owenshen@mit.edu}
\date{}

\subjclass[2020]{Primary 05B20; Secondary 11B83}
\keywords{periodic Golay pair, Golay pair, sequence compression,
Turyn multiplication}

\begin{document}

\begin{abstract}
Lumsden, Kotsireas, and Bright conjectured that the $(v/2)$-compression
of every periodic Golay pair is equivalent to one of two forms determined
by its row sums. We construct a binary Golay pair of length $100$ whose
$50$-compression is
\[
([12,2],[-4,6]),
\]
and is equivalent to neither form. Since the conjecture was established
for all $v<100$, this is a least-length counterexample. Iterated Turyn
multiplication gives counterexamples of length $10^n$ for every $n\ge2$.
\end{abstract}

\maketitle

\section{The conjecture and an obstruction}

For a sequence $X=(x_0,\ldots,x_{v-1})$, write
\[
N_X(s)=\sum_{j=0}^{v-1-s}x_jx_{j+s},
\qquad
P_X(s)=\sum_{j=0}^{v-1}x_jx_{j+s\bmod v}.
\]
A pair $(X,Y)$ of $\{\pm1\}$-sequences is a \emph{Golay pair} if
$N_X(s)+N_Y(s)=0$ for $1\le s<v$, and a \emph{periodic Golay pair}
if $P_X(s)+P_Y(s)=0$ for $1\le s<v$. Since
\[
P_X(s)=N_X(s)+N_X(v-s),
\]
every Golay pair is periodic.

For even $v$, let
\[
r(X)=\sum_{j=0}^{v-1}x_j,
\qquad
e(X)=\sum_{j=0}^{v-1}(-1)^jx_j.
\]
The $(v/2)$-compression of $X$ is the length-two sequence
\[
X^{(v/2)}
=
\left[
\frac{r(X)+e(X)}2,
\frac{r(X)-e(X)}2
\right].
\]
Lumsden, Kotsireas, and Bright conjectured that if the
$(v/2)$-compression of a periodic Golay pair has row sums $a$ and $b$,
then it is equivalent to one of
\begin{equation}\label{eq:forms}
([0,a],[0,b])
\quad\text{or}\quad
\left(
\left[\frac{a+b}{2},\frac{a-b}{2}\right],
\left[\frac{a+b}{2},\frac{b-a}{2}\right]
\right).
\end{equation}
Here equivalence is generated by swapping the sequences, cyclically
shifting or reversing the first sequence, simultaneous decimation, and
simultaneous alternating negation, as in
\cite[\S2.1 and Conjecture~1]{LKB}.
They established the conjecture for $v<100$ and identified $v=100$ as
the first unresolved length~\cite[\S7]{LKB}. These generators induce
the corresponding operations on the length-two compression: a shift or
reversal exchanges the parity classes, an admissible decimation is odd
and preserves them, and alternating negation changes the sign of the
odd class. Negation of one sequence is itself generated by the five
operations~\cite[\S2.1]{LKB}.

\begin{theorem}\label{thm:main}
For every integer $n\ge2$, there is a binary Golay pair of length
$10^n$ whose $(10^n/2)$-compression is equivalent to neither of the
two pairs in \eqref{eq:forms}, with $a,b$ equal to its row sums.
Consequently, Conjecture~1 of Lumsden, Kotsireas, and Bright is
false in infinitely many lengths. The case $n=2$ is a least-length
counterexample.
\end{theorem}

Let $\mathcal U=\{\pm1,\pm\ii\}$. For a length-two integer pair
$H=([x_0,x_1],[y_0,y_1])$, define
\[
z(H)=(x_0+x_1)+\ii(y_0+y_1),
\qquad
w(H)=(x_0-x_1)+\ii(y_0-y_1).
\]

\begin{lemma}\label{lem:obstruction}
If $H$ is equivalent to either pair in \eqref{eq:forms}, then
\begin{equation}\label{eq:unit-condition}
w(H)\in \mathcal U\,z(H)\ \cup\ \mathcal U\,\overline{z(H)}.
\end{equation}
\end{lemma}

\begin{proof}
For the two pairs in \eqref{eq:forms}, respectively, one has
$w=-z$ and $w=\ii\overline z$. At length two, swapping the sequences
sends
\[
(z,w)\longmapsto(\ii\overline z,\ii\overline w),
\]
shifting or reversing the first sequence sends
$(z,w)\mapsto(z,-\overline w)$, decimation is trivial, and simultaneous
alternating negation sends $(z,w)\mapsto(w,z)$. Since $\mathcal U$ is
closed under multiplication, inversion, and conjugation, a direct
substitution shows that each operation preserves
\eqref{eq:unit-condition}.
\end{proof}

\section{Construction}

\begin{proof}[Proof of Theorem~\ref{thm:main}]
Identify the symbols $+$ and $-$ with $1$ and $-1$. Consider the
length-$10$ pairs
\[
\begin{aligned}
U&=\texttt{+--++++++-}, &
V&=\texttt{+--+-+---+},\\
C&=\texttt{-++-+-+++-}, &
D&=\texttt{-++++++--+}.
\end{aligned}
\]
With
\[
c=(3,0,1,0,-1,-2,1,2,-1),
\]
direct calculation gives
\[
(N_U(s))_{s=1}^{9}=(N_D(s))_{s=1}^{9}=c,
\qquad
(N_V(s))_{s=1}^{9}=(N_C(s))_{s=1}^{9}=-c.
\]
Thus $(U,V)$ and $(C,D)$ are Golay pairs.

Put
\[
p_j=\frac{u_j+v_j}{2},
\qquad
q_j=\frac{u_j-v_j}{2}.
\]
Explicitly,
\[
\begin{aligned}
p&=(1,-1,-1,1,0,1,0,0,0,0),\\
q&=(0,0,0,0,1,0,1,1,1,-1).
\end{aligned}
\]
Hence $p_jq_j=0$, $p_j^2+q_j^2=1$ for every $j$, and
\begin{equation}\label{eq:sums}
\sum_{j=0}^{9}p_j=1,
\qquad
\sum_{j=0}^{9}q_j=3.
\end{equation}

For an even-length pair $(X,Y)$, let $X^{\rev}$ denote reversal and
define
\begin{equation}\label{eq:turyn}
\mathcal T(X,Y)
=
\left(
\mathop{\Vert}_{j=0}^{9}(p_jX+q_jY^{\rev}),
\mathop{\Vert}_{j=0}^{9}(-q_jX^{\rev}+p_jY)
\right),
\end{equation}
where $\Vert$ denotes concatenation. This is a standard form of
Turyn multiplication~\cite{Turyn}.

Let $F_X(z)=\sum_jx_jz^j$, and let $P,Q$ be the polynomials with
coefficient sequences $p,q$. If $(A,B)=\mathcal T(X,Y)$ and $m$ is
the common length of $X,Y$, then, using
\[
F_{X^{\rev}}(z)=z^{m-1}F_X(z^{-1}),
\]
cancellation of the cross terms gives
\begin{align*}
&F_A(z)F_A(z^{-1})+F_B(z)F_B(z^{-1})\\
&\quad=
\bigl(P(z^m)P(z^{-m})+Q(z^m)Q(z^{-m})\bigr)\\
&\qquad\quad{}\times
\bigl(F_X(z)F_X(z^{-1})+F_Y(z)F_Y(z^{-1})\bigr).
\end{align*}
Since $F_U=P+Q$ and $F_V=P-Q$, the Golay identity for $(U,V)$ yields
\[
2\bigl(P(z)P(z^{-1})+Q(z)Q(z^{-1})\bigr)=20.
\]
Thus the first factor above is identically $10$. If $(X,Y)$ is a
Golay pair of length $m$, the second factor is $2m$; therefore
\eqref{eq:turyn} is a Golay pair of length $10m$. It is binary because
$p_jq_j=0$ and $p_j^2+q_j^2=1$.

Set $(A_1,B_1)=(C,D)$ and recursively define
\[
(A_{n+1},B_{n+1})=\mathcal T(A_n,B_n).
\]
Then $(A_n,B_n)$ is a binary Golay pair of length $10^n$. Define
\[
z_n=r(A_n)+\ii r(B_n),
\qquad
w_n=e(A_n)+\ii e(B_n).
\]
If $H_n$ is its $(10^n/2)$-compression, then
$z(H_n)=z_n$ and $w(H_n)=w_n$.

Since the block length in \eqref{eq:turyn} is even, alternating
sums add across the concatenation; moreover,
$e(X^{\rev})=-e(X)$ for every even-length sequence $X$. Using
\eqref{eq:sums} in \eqref{eq:turyn} therefore gives
\[
\begin{aligned}
r(A_{n+1})&=r(A_n)+3r(B_n), &
r(B_{n+1})&=-3r(A_n)+r(B_n),\\
e(A_{n+1})&=e(A_n)-3e(B_n), &
e(B_{n+1})&=3e(A_n)+e(B_n).
\end{aligned}
\]
Consequently,
\begin{equation}\label{eq:recurrences}
z_{n+1}=(1-3\ii)z_n,
\qquad
w_{n+1}=(1+3\ii)w_n.
\end{equation}
The initial values are
\[
z_1=2+4\ii,
\qquad
w_1=4-2\ii=-\ii z_1.
\]
Writing $t=1+3\ii$, equations \eqref{eq:recurrences} give
\begin{equation}\label{eq:closed}
z_n=\overline t^{\,n-1}z_1,
\qquad
w_n=-\ii t^{n-1}z_1.
\end{equation}

Suppose $n\ge2$ and \eqref{eq:unit-condition} holds. If
$w_n=uz_n$ for some $u\in\mathcal U$, then \eqref{eq:closed} implies
that $t^{n-1}$ and $\overline t^{\,n-1}$ are associates in
$\mathbb Z[\ii]$. This is impossible, since
\[
t=(1+\ii)(2+\ii),
\qquad
\overline t=(1-\ii)(2-\ii),
\]
and $2+\ii$ and $2-\ii$ are nonassociate Gaussian primes, both of
norm $5$. Unique factorization in $\mathbb Z[\ii]$ rules out the
association.

If instead $w_n=u\overline{z_n}$, cancellation in
\eqref{eq:closed} gives $-\ii z_1=u\overline{z_1}$, but
\[
\frac{-\ii z_1}{\overline{z_1}}
=
\frac{4+3\ii}{5}
\notin\mathcal U.
\]
Thus \eqref{eq:unit-condition} fails, and
Lemma~\ref{lem:obstruction} proves the asserted inequivalence.

For $n=2$, equations \eqref{eq:closed} give
\[
z_2=14-2\ii,
\qquad
w_2=10+10\ii,
\]
so the compression of $(A_2,B_2)$ is
\[
H_2=([12,2],[4,-6]).
\]
Replacing $B_2$ by $-B_2$ gives another Golay pair whose compression is
\[
\widetilde H_2=([12,2],[-4,6])
\]
with row sums $a=14$ and $b=2$. The two representatives in
\eqref{eq:forms} are then
\[
([0,14],[0,2])
\quad\text{and}\quad
([8,6],[8,-6]).
\]
The maximum absolute entry is respectively $12$, $14$, and $8$;
this quantity is invariant under the length-two equivalence
operations, giving a direct check at length $100$. Since the
conjecture holds for $v<100$ by~\cite{LKB}, this counterexample has
least possible length.
\end{proof}

\begin{thebibliography}{9}

\bibitem{LKB}
T.~Lumsden, I.~Kotsireas, and C.~Bright,
\emph{New results on periodic Golay pairs},
Math. Comp. \textbf{95} (2026), no.~359, 1517--1539,
\href{https://doi.org/10.1090/mcom/4096}
{doi:10.1090/mcom/4096}.

\bibitem{Turyn}
R.~J. Turyn,
\emph{Hadamard matrices, Baumert--Hall units, four-symbol sequences,
pulse compression, and surface wave encodings},
J. Combin. Theory Ser. A \textbf{16} (1974), no.~3, 313--333,
\href{https://doi.org/10.1016/0097-3165(74)90056-9}
{doi:10.1016/0097-3165(74)90056-9}.

\end{thebibliography}

\end{document}